% 第9章 图神经网络
\section{图神经网络：关系结构与跨资产信息}\label{sec:gnn}

前面几族大多把每只资产当成独立的样本：给资产 $i$ 一组特征，预测资产 $i$ 的收益。但市场的现实是资产之间相互牵连——供应商的坏消息会传到客户，同行业龙头的估值会带动板块，共同股东的流动性冲击会同时打中好几只票。图神经网络（GNN）的价值，就是把这层“谁和谁有关系”的结构显式地建进模型。

\subsection{问题设定}

把市场建成图 $\mathcal{G}=(\mathcal{V},\mathcal{E})$：节点 $\mathcal{V}$ 是资产/公司，边 $\mathcal{E}$ 是某种经济关系（行业、供应链、共同股东、新闻共现、价格相关）。GNN 通过消息传递更新节点表示：
\begin{equation}
  h_i^{(l+1)} = \sigma\Big( \sum_{j\in\mathcal{N}(i)} \alpha_{ij}\, W\, h_j^{(l)} \Big),
\end{equation}
$\alpha_{ij}$ 是边权（注意力，如 GAT\cite{velickovic2018gat}）。对交易而言，“图”提供了横截面之外的\textbf{结构性先验}：它规定了信息沿哪些通道传导，从而在低信噪比下减少搜索空间。

\subsection{一条从“关系”到“可交易”的证据线}

\textbf{经济关系的可预测性（前 GNN 时代）}。Cohen 与 Frazzini（2008）\cite{cohenfrazzini2008} 发现供应链上的动量会从供应商传导到客户，是“关系 $\to$ 可交易收益”的经典证据。这条证据很重要，因为它说明图上的边不只是建模便利，而是有真实的经济内容。

\textbf{GNN 进入股票预测}。Feng 等（2019）\cite{feng2019rsr} 的时序关系排序模型把公司间的时序关系建进排序学习；Chen、Wei 与 Huang（2018）\cite{chen2018gcn} 用图卷积引入公司关系；Matsunaga、Suzumura 与 Takahashi（2019）\cite{matsunaga2019} 系统检验了 GNN 在股票预测上的表现，是较早的严格对照研究之一。这些工作确认了一个结论：\textbf{关系结构确实带来增量，但增量的大小高度依赖于“边是怎么定义的”。}

\textbf{注意力与异质图}。Xu 等（2021）\cite{xu2021hats} 的 HATS 用层次化图注意力建模个股与概念/行业的多层关系；Wang 等（2021）\cite{wang2021thgnn} 用时序异质图同时建模多种关系；Sawhney 等（2020）\cite{sawhney2020} 把社交媒体文本与公司相关图结合做股价动向预测。这一代的趋势是把“单一关系图”升级成“多关系、时变的异质图”，逼近真实市场的复杂度。

\subsection{边怎么定义：五种做法的取舍}

这一族最实际、也最少被系统讨论的问题，是“边从哪来”。不同的边构造方式，直接决定了模型能捕捉什么、以及会不会制造假增量。

\begin{table}[htbp]
\centering
\caption{常见边构造方式的取舍}
\label{tab:edges}
\begin{tabular}{@{}>{\raggedright\arraybackslash}p{2.6cm}>{\raggedright\arraybackslash}p{4.2cm}>{\raggedright\arraybackslash}p{3.4cm}>{\raggedright\arraybackslash}p{4.0cm}@{}}
\toprule
边类型 & 构造方式 & 优点 & 主要风险 \\
\midrule
行业/板块 & 官方分类标签 & 稳定、全市场可观测 & 粒度粗，信息多已被价格反映 \\
供应链 & 供应商--客户披露 & 经济含义强、传导机制清楚 & 披露稀疏、滞后，覆盖有限 \\
共同股东/机构持仓 & 13F 等持仓数据 & 反映资金联动 & 财报/持仓有季度级滞后 \\
价格相关 & 滚动相关系数 & 覆盖全、易获得 & 信息已在价内，易产生假增量 \\
新闻/文本共现 & 语义相似度建边 & 覆盖广、更新快 & 噪声大，时点对齐困难 \\
\bottomrule
\end{tabular}
\end{table}

取舍原则很朴素：\textbf{越有经济含义的边越可信，但越稀疏、越滞后；越容易构造的边越有“假增量”风险，因为价格相关的信息大多已经被价格吸收。}因此一个务实的做法是“多种边并用 + 分别做消融”，确认增量到底来自哪一种边——如果去掉价格相关边后增量消失，那这个“图”的贡献就值得怀疑。

\subsection{一个可执行的关系消融协议}

“图有没有用”这个问题，光靠报告一个更好的指标回答不了，需要一套消融设计。建议按下面四步做，并用同一口径（同数据、同成本、同评测协议）比较：

\begin{enumerate}
  \item \textbf{基线}：与目标模型同架构、但不含任何图结构（把邻居信息去掉）。这是判断“图是否带来增量”的唯一参照。
  \item \textbf{只加经济含义边}：仅用行业分类、供应链披露、共同股东这类有经济内容的边重建模型。
  \item \textbf{再加价格相关边}：在第 2 步基础上加入滚动相关边，看是否有进一步提升。
  \item \textbf{安慰剂检验}：把边随机重连（保持度分布）后重跑。若“随机图”也带来同样提升，说明所谓的增量来自模型容量而非关系信息。
\end{enumerate}

判定规则可以直接写死：\textbf{只有当第 2 步显著优于第 1 步、且第 4 步不显著优于第 1 步时，图结构的贡献才算成立。}如果增量主要来自第 3 步的价格相关边，则应如实报告“增量来自价格相关性的再编码”，而不是“关系信息带来了新 alpha”——这两句话在交易上含义完全不同，前者几乎必然在扣成本后消失。

\subsection{落地要素}

\begin{itemize}
  \item \textbf{边的时点必须可观测}。用“今天的行业分类”去回测十年前，是典型的前瞻偏差——公司的主营、供应链关系都在变。行业分类、供应链披露、共同股东名单，都必须用当时已知的版本。
  \item \textbf{边不能太“软”}。用价格相关性构图最省事，但价格相关本身已被价格反映，用它会制造“信息已在价内”的假增量。有经济含义的边（供应链、共同股东）更可信也更难得。
  \item \textbf{图规模与计算成本}。全市场图的消息传递计算量随节点数增长，日频更新在工程上需要权衡。
\end{itemize}

\subsection{复现状态与失效边界}

这一族多数工作在公开数据集（如基于新闻共现的公司关系集）上评测，但\textbf{数据集的时间划分、是否含退市股、关系是否 point-in-time，往往披露不足}，导致“同一方法不同论文结论相反”。因此复现状态整体为\textbf{部分复现}：架构可复现，结论不可直接比较。

失效边界有两条比较特殊。第一，\textbf{关系被价格吸收}：如果某种关系早已被市场定价，则把它写进图不产生新信息，只增加过拟合自由度。第二，\textbf{图是人为选定的}：你决定建哪些边，就等于先验地决定了信息沿哪些通道流动——\textbf{把假设藏进图结构，是这一族最容易被忽视的自欺。}本章要记住的一句话是：\textbf{图的价值在于把有经济含义的关系显式写进模型，而它的风险也恰在于——你写进去的关系，可能只是你自己的假设。}
